\chapter{Fundamentals and Methods of Polarized Resonant Soft X-ray Reflectivity}\label{chap:methods}

This chapter develops the physical and practical framework used in every later chapter. It first describes what a far-field scattering measurement records and why specular reflectivity, unlike transmission scattering, must be computed beyond the Born approximation. It then introduces the gradient model used for interfacial structure in Chapter 4, built on the film tensor of Chapter 2. The remaining sections describe how polarized resonant reflectivity is measured at ALS beamline 11.0.1.2, from alignment through reduction, stitching, and uncertainty, and how it is analyzed, from the forward model through optimization, posterior estimation, and model selection.

\section{\texorpdfstring{Scattering in the Far Field}{Scattering in the Far Field}}\label{figmarker9x-scattering-in-the-far-field}

\begin{figure}[tbp]
  \centering
  \includegraphics[width=\linewidth,height=0.75\textheight,keepaspectratio]{fig2-01_scattering-geometry}
  \caption[Scattering geometries and the momentum transfer.]{Scattering geometries and the momentum transfer. a) A plane wave with wavevector k scatters from a finite object, and the curved near-field wavefronts flatten into the far-field pattern recorded along k\ensuremath{'} at scattering angle 2θ. The momentum transfer is q = k\ensuremath{'} − k, with |q| = (4π/λ) sin θ. b) Specular reflection from a film on a substrate at grazing angle θ, for which qz = (4π/λ) sin θ lies along the film normal. s polarization places E normal to the plane of incidence and p polarization places E in it.}
  \label{fig:3.1}
\end{figure}


Every measurement in this work records intensity in the far field as a function of the momentum transfer \(\mathbf{q} = \mathbf{k}' - \mathbf{k}\) between the incident and scattered wavevectors (Fig.~\ref{fig:3.1}). For elastic scattering \(\left| \mathbf{k} \right| = \left| \mathbf{k}' \right| = 2\pi/\lambda\), so the magnitude \(\left| \mathbf{q} \right| = (4\pi/\lambda)\sin\theta\) is fixed by the scattering angle \(2\theta\) alone (Fig.~\ref{fig:3.1}a).\autocite{attwood1999soft} In transmission RSoXS, \(\mathbf{q}\) is read from the position of the diffuse intensity on an area detector, with the pattern reporting in-plane structure through its dependence on \(\left| \mathbf{q} \right|\).\autocite{collins2012polarized,gann2016origins,grabner2025resonant} Specular reflectivity instead holds the detector at twice the grazing angle, so the reflected beam leaves at the angle of incidence and \(\mathbf{q} = q_{z}\widehat{\mathbf{z}}\) lies along the film normal (Fig.~\ref{fig:3.1}b). To first order, scanning \(\theta\) then samples the Fourier components of the depth profile of the optical constants, which is why reflectivity resolves stratification that transmission averages over.\autocite{wang2005soft,wang2007resonant} The plane of incidence also defines the two polarizations used throughout. In s polarization \(\mathbf{E}\) lies normal to the plane of incidence and therefore in the film plane at every \(\theta\). In p polarization \(\mathbf{E}\) lies in the plane of incidence and carries a component \(\cos\theta\) along the film normal, largest at grazing incidence. For a uniaxial film the two polarizations therefore weight the in-plane and out-of-plane components of the optical tensor differently, which is the source of the orientation sensitivity used in Chapters 4 and 5.\autocite{mezger2011molecular,thelen2020molecular}

The far-field condition is stated through the Fresnel number \(F = a^{2}/(L\lambda)\), which compares the size \(a\) of a scattering feature with the distance \(L\) to the detector and the wavelength \(\lambda\).\autocite{foldy1945multiple} At 284 eV, \(\lambda \approx 4.37\) nm, and for \(a = 10\) nm domains with \(L \approx 0.5\) m, \(F \approx 5 \times 10^{- 8}\). \(F\) approaches unity only for structures near \((L\lambda)^{1/2} \approx 47\) µm, far larger than any feature studied here, so every measurement in this work is a Fraunhofer measurement. This condition concerns the detector and says nothing about the field inside the sample. Figure~\ref{fig:3.2} follows a plane wave through a small cluster of scatterers with a two-dimensional multiple-scattering calculation. Near the cluster each scatterer is driven by the incident wave and by the fields scattered from its neighbors, so the crests of the total field are distorted\autocite{lax1951multiple,foldy1945multiple}, with a scatterer deep in the cluster radiating in response to both. With distance these distortions propagate out into smooth outgoing wavefronts, and the detector records only the far-field intensity \(I(2\theta)\), a forward peak set by the spacing and size of the scatterers with weaker side maxima. The measurement therefore discards the phase of the far field and the structure of the near field, and the sample must be inferred by forward modeling and fitting, as described in the section on fitting. When the interaction between scatterers is weak, one scattering event per photon suffices and the first Born approximation \(A\left( \mathbf{q} \right) = \int\Delta n\left( \mathbf{r} \right)\, e^{i\mathbf{q} \cdot \mathbf{r}}\, d^{3}r\) applies, which is the basis of transmission RSoXS and of its simulation.\autocite{collins2022resonant,gann2016origins,saurabh2023textless,dudenas2025how,grabner2024molecular} When it is strong, as for specular reflection near the critical angle, the multiple scattering of Fig.~\ref{fig:3.2} must be solved exactly, which the transfer matrix method of Fig.~\ref{fig:3.6} does for a layered film.

\begin{figure}[tbp]
  \centering
  \includegraphics[width=\linewidth,height=0.75\textheight,keepaspectratio]{fig2-02_near-far-field}
  \caption[The far-field pattern emerges from local interactions between scatterers.]{The far-field pattern emerges from local interactions between scatterers. A plane wave with wavevector k and intensity I0 illuminates five scatterers of size a separated by distances of order d. The front scatterer scatters forward along ks, so a deeper scatterer is driven by both I0 and this secondary field before it radiates along k\ensuremath{'}. Black lines are crests of the total scattered field from a two-dimensional multiple-scattering (Foldy-Lax) calculation, distorted near the cluster and flattening into outgoing wavefronts with distance. A detector at L ≫ d, a records only the far-field intensity I(2θ) (red), and the L axis is broken.}
  \label{fig:3.2}
\end{figure}


\section{Why Specular Reflectivity Is Not Kinematic}\label{why-specular-reflectivity-is-not-kinematic}

Specular reflectivity is the measurement in this work for which the first Born approximation fails. Below the critical angle \(\theta_{c} = (2\delta)^{1/2}\) reflection is nearly total, and near \(\theta_{c}\) the reflected amplitude is comparable to the incident amplitude, so the field inside the film differs from the incident field and waves reflected between interfaces interfere with one another.\autocite{parratt1954surface} Far above \(\theta_{c}\) the kinematic result \(R\left( q_{z} \right) \approx R_{F}\left( q_{z} \right)\left| \delta_{\infty}^{- 1}\int(d\delta/dz)\, e^{iq_{z}z}\, dz \right|^{2}\) holds, where \(R_{F}\) is the Fresnel reflectivity of the bare substrate, \(\delta(z)\) the depth profile of the refractive decrement, and \(\delta_{\infty}\) its value in the substrate.\autocite{alsnielsen2011elements} It is used here only to read two features of the data. \(R_{F}\) decays as \(q_{z}^{- 4}\), which the \(Rq_{z}^{4}\) representation used in later chapters removes, and a layer of thickness \(t\) produces Kiessig fringes spaced by \(\Delta q_{z} \approx 2\pi/t\), so the fringe period reports the thickness directly.\autocite{kiessig1931interferenz} Near \(\theta_{c}\), where the data are most intense, the reflectivity is instead computed exactly in one dimension by the transfer matrix method described in the section on simulating reflectivity, so no perturbative approximation enters the fits in this work.

\section{Orientation gradient theoretical foundation}\label{orientation-gradient-theoretical-foundation}

We formulate our gradient model from a Ginzburg-Landau framework of second order phase transitions, where we treat the molecular orientation and density as order parameters (not to be confused with nematic order parameters \(S(\gamma)\) or \(P(\gamma)\)).\autocite{gennes2023physics,ginzburg2009theory} Along the one dimensional depth axis \(z\) of a film of thickness \(L\), we take the order parameters as the deviation in the density \(\Delta\rho(z) = \rho(z) - \rho_{b}\) and orientation \(\Delta\gamma(z) = \gamma(z) - \gamma_{b}(z)\) from the bulk. The total free energy of the film is described by the functional

\[\begin{array}{r}
F\lbrack\Delta\rho,\Delta\gamma\rbrack = \int_{0}^{L}{f(\Delta\rho,\Delta\gamma) + \frac{1}{2}c_{\rho}\left( {\Delta\rho}' \right)^{2} + \frac{1}{2}c_{\gamma}\left( {\Delta\gamma}' \right)^{2}dz} \\
 + F_{v}\left\lbrack \Delta\rho(0),\Delta\gamma(0) \right\rbrack + F_{Si}\left\lbrack \Delta\rho(L),\Delta\gamma(L) \right\rbrack
\end{array}\]

Here \(f(\Delta\rho,\Delta\gamma)\) is the local bulk free energy density, and the terms proportional to \(c_{\rho}\) and \(c_{\gamma}\) are the gradient energy penalties that oppose abrupt variations within the film. The interface terms \(F_{v}\) and \(F_{Si}\) account for the anchoring potentials at the interface boundaries. Expanding \(f(\Delta\rho,\Delta\gamma)\) to quadratic order about the bulk minimum,

\[\begin{array}{r}
f(\Delta\rho,\Delta\gamma) \approx f_{0} + \frac{1}{2}a_{\gamma}{\Delta\gamma}^{2} + \frac{1}{2}a_{\rho}{\Delta\rho}^{2} + b\Delta\rho\Delta\gamma.
\end{array}\]

where \(a_{\rho}\) and \(a_{\gamma}\) are the stability coefficients of the bulk phase, and \(b\) the bilinear coupling term between the density and orientation fields.

Minimizing the total variation \(\delta F = 0\) and integrating the gradient terms by parts yields a system of coupled second-order differential equations for the film interior,

\[{\begin{array}{r}
c_{\rho}{\Delta\rho}^{''} = a_{\rho}\Delta\rho + b\Delta\gamma\ 
\end{array}
}\begin{array}{r}
c_{\gamma}{\Delta\gamma}^{''} = a_{\gamma}\Delta\gamma + b\Delta\rho\ 
\end{array}\]

with boundaries dictated by the surface free energies. To leading order in the interface fields \(h\), the boundary variation yields Robin boundary conditions, where the gradient of the order parameter is proportional to the surface field, i.e.

\[\begin{array}{rlr}
c_{\rho}\Delta\rho'(0) = - h_{\rho}^{v}\  & & c_{\rho}\Delta\rho'(L) = h_{\rho}^{Si}
\end{array}\]

and with similar form for \(\Delta\gamma\). When the surface fields \(h\) are much stronger than the bulk stiffness coefficients \(c/\tau\), the system enters the Dirichlet regime of strong anchoring, in which the interface potentials set the order parameters to fixed interfacial values, \(\Delta\rho(0) = {\Delta\rho}_{v}\) and \(\Delta\rho(L) = {\Delta\rho}_{Si}\). Within this regime the interfaces impose an orientational state regardless of the bulk interaction, a feasible state given the known different growth modes for ZnPc films at the substrate and surface.\autocite{derrien2022interfacial,yang2015molecular}

This system of equations may be decoupled into a normal modes problem, under the coordinate transform \(\widetilde{\Delta\rho} = \sqrt{c_{\rho}}\Delta\rho\) and similar transform for \(\Delta\gamma\). Given that the bulk represents the stable minimum, the general solution is restricted to the linear combination of exponentials decaying from the boundaries with coherence lengths \(\tau_{\pm}\) derived from the eigenvalues of the coupled system. When the orientation and density coupling is nonzero, these coherence lengths are shared between the order parameters. As each interface bears its unique interactions that dictate this coherence length, we assume that each interface is dominated by a single mode. The interface gradient model used in the main text satisfies these conditions,

\[{\begin{array}{r}
\Delta\rho(z) = \Delta\rho_{v}\exp\left( - z/\tau_{vac} \right) + \Delta\rho_{Si}\exp\left( - (L - z)/\tau_{Si} \right)
\end{array}
}\begin{array}{r}
\Delta\gamma(z) = \Delta\gamma_{v}\exp\left( - z/\tau_{vac} \right) + \Delta\gamma_{Si}\exp\left( - (L - z)/\tau_{Si} \right)\ 
\end{array}\]

While this model is congruent with the continuous reorientation under an evolving growth front, it remains an empirical approximation of the underlying interfacial physics. A more rigorous determination of the anchoring potentials would require a detailed investigation of the boundary conditions, specifically whether the system is governed by Dirichlet or Robin boundaries as this determines the relationship between exponential amplitudes and the interfacial potentials. Furthermore, this model assumes a continuous structural relaxation; however, it is possible that the molecular realignment across the interface occurs under a different mechanism than described here requiring a different theoretical basis.

\section{Polarized Resonant Reflectivity}\label{polarized-resonant-reflectivity}

Polarized resonant reflectivity measures the film tensor of the preceding sections as a function of depth. It is the only technique with demonstrated sensitivity to molecular orientation at buried heterointerfaces.\autocite{ferron2023composition,heilman2026sub} Orientation at such interfaces sets the performance of organic photovoltaic devices\autocite{tumbleston2014influence} and the efficiency of organic light-emitting diodes,\autocite{yokoyama2011molecular} yet few studies have resolved the depth profiles of orientation and composition together. The limitation lies in the analysis more than in the measurement, since existing frameworks were built for single systems and carry assumptions that fail elsewhere. The sections that follow describe the measurement and a general analysis framework for it. Every program and algorithm is open source and maintained jointly by the Collins group at WSU, the ALS computing group, and the beamline scientists at 11.0.1.2, where the analysis is being integrated into the beamline control system toward autonomous experimentation. The framework is demonstrated here for reflectivity but applies to RSoXS measurements generally.

\section{Experimental Methods}\label{experimental-methods}

Reflectivity complements the transmission geometries above by directly resolving depth, mapping a film\textquotesingle s optical constants as a function of distance from the surface rather than averaging over the full film thickness. Polarization further resolves the orientational anisotropy captured in the tensor framework above, since \emph{s} and \emph{p} polarized light couple to different components of the optical tensor at a given angle of incidence. Together, depth and polarization sensitivity make reflectivity the natural geometry for testing the orientation and chemical-state profiles this dissertation is built around.

X-ray reflectivity measures the specular reflected intensity, directly probing the reflection coefficients set by the refractive index of each layer. In multi-layered films, multiple reflections at each interface break the Born approximation requiring a rigorous simulation platform that can express these multiple reflection and transmission events. For anisotropic materials, the standard approach follows the transfer matrix method that recursively solves the reflection and transmission coefficients at the interfaces of thin slabs bearing uniform optical constants. This means that not only is the reflected intensity set by the surface optical constants, but also those buried in every layer below the surface, allowing the technique to resolve optical stratification that no other technique can. The final product of the technique is a depth profile of the optical constants along the depth of a film, and in polarized and resonant reflectivity, interpreting this profile is core to resolve molecular state as a function of depth.

Real interfaces are rough or graded, which damps the specular reflectivity at each interface and redistributes intensity into diffuse scattering. For organic films the distinction matters, since a sharp but rough interface and a gradual change of composition or orientation can damp the specular signal similarly while implying different physics at the interface. The interface models used in this work, and the regimes in which each applies, are described in the section on simulating reflectivity (Fig.~\ref{fig:3.7}).

\subsection{\texorpdfstring{Reflectivity setup at 11.0.1.2}{Reflectivity setup at 11.0.1.2}}\label{figmarker11x-reflectivity-setup-at-11.0.1.2}

\begin{figure}[tbp]
  \centering
  \includegraphics[width=\linewidth,height=0.75\textheight,keepaspectratio]{beamline-schematic}
  \caption[Beamline schematic of ALS Beamline 11.0.1.2.]{Beamline schematic of ALS Beamline 11.0.1.2. a) Optical path of X-rays: (1) UPS undulator injection source producing polarized photons, (2) toroidal focusing mirror, (3) beam shaping slits, (4) vertical deflection steering and focusing mirrors, (5) grating monochromator for energy selection, (6) higher order suppression optics to remove higher harmonics and linearly attenuate the beam, (7) gold mesh intensity monitor, (8) reflectometer chamber. b) Interior of the 11.0.1.2 chamber: (1) beam entrance path after leaving the beam shaping slits, (2) sample stage with four rotation axes, (3) area detector with three degrees of freedom mounted on a goniometer rotation arm with near coaxial rotation with the sample stage.}
  \label{fig:3.3}
\end{figure}


Resonant reflectivity is recorded on the area detector (Fig.~\ref{fig:3.3}b, 3) after reflection from the sample (Fig.~\ref{fig:3.3}b, 2), with each measurement following four steps.

\begin{enumerate}
\def\labelenumi{\arabic{enumi}.}
\item
  \textbf{Beam Shaping} -- The beam defining slits, higher order suppressor, and focusing mirrors are fine tuned to block anomalous scattering from the optical elements while lowering the intensity so the direct beam can safely be measured on the area detector without damage.
\item
  \textbf{Sample Alignment} -- The sample theta and height (Fig.~\ref{fig:3.3}b, \(\theta_{s}\) and \(z\) respectively) are iteratively aligned to the direct beam so that the \(\theta_{s} = 0\) corresponds to perfectly grazing, with \(z\) set to cut the beam intensity in half.
\item
  \textbf{Scan Definition} -- The experimental orchestration scan is configured to capture upwards of 10 orders of magnitude of signal as a function of sample and detector angles \(\theta_{s\ }\ \&\ \theta_{d}\) while maintaining normalizability.
\item
  \textbf{Scan Running and Monitoring} -- Finally, scans are monitored over the angular range watching for frames where the reflected intensity drops below the detector dynamic range, or for instances where the reflected intensity saturates the detector.
\end{enumerate}

This general procedure is repeated for each measurement condition (photon energy, photon polarization, etc.) to fully capture the samples reflectivity coefficients over these conditions. The exact procedures for this technique are evolving, with recent developments towards fully agentic data collection (Fig.~\ref{fig:3.12}).

\subsection{Sample Alignment}\label{sample-alignment}

Reflectivity normalization is performed by iteratively refining the sample height (sample \(z\) motor) and the sample \(\theta_{s}\) motors such that the beam is perfectly flat and cutting the beam in half. This is performed by exploiting the size of the beam itself. First using the nominal z and theta positions, the sample is rotated flat and below the beam. From there it is incrementally moved through the beam and the intensity is watched on the photodiode. This results in a sigmoid distribution that convolves the beam profile with edge of the sample and sample stage, a process known as a knife edge test. The place where the beam intensity is halved centers the topmost point of the sample into the center of the beam.

Next the \(\theta_{s}\) motor is scanned around zero degrees, \(\theta_{s} \in ( - 1{^\circ},1{^\circ})\). When the sample is poorly aligned, the photodiode bears a wide distribution centered around the point where the intensity is maximized. The sample \(\theta\) motor is then reset so that this angle is taken as \(0{^\circ}\) before returning to the

sample \(z\) alignment routine. As the alignment improves these distributions change shape, the sample z becomes a sharper distribution as it is no longer possible to reflect some of the incident beam as the sample moves into position. Yet the sample theta positions change more drastically.

\begin{figure}[tbp]
  \centering
  \includegraphics[width=\linewidth,height=0.75\textheight,keepaspectratio]{fig2-07_alignment}
  \caption[Sample alignment with a finite beam and a finite photodiode, simulated by ray Monte Carlo.]{Sample alignment with a finite beam and a finite photodiode, simulated by ray Monte Carlo. a, b) Knife-edge geometry, in which the sample is raised through the beam at θ = 0, and the photodiode signal in 25 µm steps with an error-function fit, which places the surface at z0 = −0.7 ± 0.5 µm for a beam of σ ≈ 41 µm. c, d) Rocking geometry at fixed z0 with the detector at 0°, where light reflected from the surface (dark red) adds to the unblocked beam until a larger tilt sends it over the diode (gray, dashed), and the scan in 0.05° steps with a forward-model fit (red). The resulting cliff near θs = 0.29° shifts the half-maximum midpoint (gray, dashed) to 0.06° for a true offset of zero, while the fit (red, dashed) recovers an offset of 0.001 ± 0.003°. e, f) θ-2θ geometry with the detector at θd = 2°, where the whole beam reflects onto the diode near θs = 1°, and the scan in 0.04° steps with a soft-edged top-hat fit centered at θ0 = 1.003 ± 0.001°. Angles in the schematics are exaggerated.}
  \label{fig:3.4}
\end{figure}


Due to the macroscopic beam size, a new phenomenon is captured on the photodiode. As theta is rotated a portion of the beam is unblocked by the sample so the intensity increases until the point where the sample is flat. But at this stage roughly 50\% of the beam is still blocked by the stage. After this turning point, the remaining beam starts to be reflected by the sample and captured on the finite sized photodiode, leading to an increase in intensity after the angle that should be considered flat until the reflected portion of the beam falls off the photodiode, leading to a sharp cliff in the photodiode signal. In this regime the half-maximum midpoint is biased toward the cliff, by 0.06° for the geometry of Fig.~\ref{fig:3.4}, so the center is taken instead from a forward-model fit to the full rocking curve (Fig.~\ref{fig:3.4}d). Another scheme takes the detector angle to \(2{^\circ}\) and the sample theta motor is then scanned over the range \(\theta_{s} \in (0{^\circ},2{^\circ})\). This approach places the entire beam reflected onto the photodiode detector through the entirety of the angular alignment process. However, the added motor movement to tilt the sample and detector can be a significant time commitment, so the best alignments instead swap to this more time-consuming mode after the sample has been roughly aligned in the flat geometry.

Figure~\ref{fig:3.4} simulates the three alignment scans by ray Monte Carlo for a beam of \(\sigma = 40\) µm with 1 mrad divergence, a 10 mm sample, and a 3 mm photodiode 150 mm downstream, to show where each scan places its center. In the knife-edge scan the sample is raised through the beam at \(\theta = 0\) (Fig.~\ref{fig:3.4}a), with the photodiode signal falling as an error function whose midpoint marks the surface and whose width returns the beam size (Fig.~\ref{fig:3.4}b). The fit places the surface at \(z_{0} = - 0.7 \pm 0.5\) µm, well within one 25 µm step. The rocking scan at fixed \(z_{0}\) is harder to read (Fig.~\ref{fig:3.4}c). Tilting the sample unblocks the lower part of the beam and also reflects the part that strikes the surface onto the diode, so the signal keeps rising past the flat angle until the reflection passes over the top of the diode and the signal drops sharply (Fig.~\ref{fig:3.4}d). The cliff makes the peak asymmetric, with the half-maximum midpoint landing at 0.06° for a sample that is flat. At 285 eV this offset shifts every angle of a scan by the same amount, a 6 percent error in \(q_{z}\) at \(\theta\) = 1°. The bias comes from the half-maximum estimator, which assumes a symmetric peak. A forward model of the same geometry, which integrates the beam profile and divergence over the sample and the diode aperture, fits the asymmetric curve directly with the amplitude, the angular offset, the beam size, and the diode acceptance free. It recovers an offset of 0.001 ± 0.003° and a beam of \(\sigma = 40.0 \pm 0.3\) µm (Fig.~\ref{fig:3.4}d, red). Fixing the diode acceptance at its true value reduces the uncertainty to 0.0013°, but an acceptance wrong by 10 percent then biases the offset by about 0.02° and more than doubles \(\chi^{2}\), so the acceptance is better left free. In the \(\theta\)-\(2\theta\) scan the detector is held at \(\theta_{d}\) = 2° and the whole beam reflects onto the diode as \(\theta_{s}\) passes through 1° (Fig.~\ref{fig:3.4}e). The signal is then a flat-topped peak with two edges set by the same geometry, and a soft-edged top-hat fit locates its center at \(\theta_{0}\) = 1.003 ± 0.001° (Fig.~\ref{fig:3.4}f), about 33 times finer than the 0.04° step and without any model of the diode. Both fits reach sub-step precision because they use the known shape of the signal rather than the position of any single point, as in localization microscopy,\autocite{thompson2002sampling} and neither improves the resolution of two neighboring features, which remains set by the step. A model fit of the rocking scan therefore sets the angular zero in the flat geometry to a few thousandths of a degree, with the slower \(\theta\)-\(2\theta\) scan providing an independent check that does not depend on the diode geometry.

The sample should be realigned after every change of sample and every change of polarization, to account for different mounting conditions between samples, and optics driven beam drifts. Critically when the polarization changes, the heat load on the optics changes, so at least 15 min should be given for the heat load to change on the optics so the alignment can be on the steady state beam location. Additionally, the beam can drift between energies due to peculiarities in the monochromator, resonant contaminants on the mirrors, as well as sample motor drift, so the beam alignment should be monitored. A mis-aligned beam will drift noticeably over the area detector to the point that it falls off when poorly aligned.

\subsection{Scan Configuration}\label{scan-configuration}

Most commonly, high-quality reflectivity data is taken as a function of incident angle. The angle-dependent reflectance can stretch over as much as 10 orders of magnitude, outpacing the dynamic range of the best area detectors. The pixel well depth of the area soft X-ray detectors is \(2^{16} = 65,536\) , four orders of magnitude. Yet thermal and electrical noise present as a noise floor that reduces that dynamic range. In the old CCD detector used through most of the instrument's life cycle this noise floor was around 500 counts, while the new CMOS detector has a floor closer to 30 leading to an increase in the dynamic range. Regardless this dynamic range requires the reflectivity profile to be captured in chunks or series that must then be combined to produce the final reflectivity curve over the full q-range.

To optimally select the angle ranges belonging to each series, the reflectivity profile needs to be known ahead of time. The best approach is to coarsely and manually scan over the angular range of the experiment and identify angles where the beam is lost below the noise floor. Then the task is to increase the intensity of the incident beam within each series to maximize the reflected intensity as measured on the detector. This is achieved through two main optical implements. First the higher order suppressor can attenuate around three orders of magnitude of intensity, so this motor is the first to control allotting for around half the usual measurement dynamic range. The second is the exposure time, which scales the recorded intensity multiplicatively. With exposures from the shutter limit, 0.001 s for the original shutter and 0.002 s for its replacement, to 10 s, this unlocks around three more orders of magnitude. In total these motors allot for nine orders of magnitude of dynamic range, just short of the sometimes required ten. The remaining range is achieved through clever sub noise floor counting, but in some cases, this precludes the observation of deep fringes and other features requiring extreme dynamic ranges of sensing.

\subsection{Reflectivity reduction and stitching}\label{reflectivity-reduction-and-stitching}

Directly analyzing the 2D reflectivity images is not feasible, so the data must be processed into a series of reflectivity curves \(I(Q)\) for each energy \(E\) and polarization. This requires extracting the reflected intensity as a fraction of the incident beam intensity from the 2D image. This algorithm has evolved continuously to accommodate new signal processing techniques, but the general approach is as follows.

\begin{enumerate}
\def\labelenumi{\arabic{enumi}.}
\item
  \textbf{Edge Correction} -- Edge effects and nonlinearities at the edges of the detector manifest in higher intensities than the rest of the flat detector. To correct for this, a border of 5 pixels is sliced from the edges within each frame.
\item
  \textbf{Row-by-row Correction} -- Running down each row, the first or last 10 pixels have their intensities averaged and that average subtracted down the row. This accounts for row dependent noise that can build up based on detector machinery. In the CCD detector this was more noticeable, but for the CMOS it is less critical. For the CMOS detector, it is also a good idea to perform a column-by-column subtraction on top of this. The result is the flattened background away from the reflected beam that removes any systemic readout noise.
\item
  \textbf{Beam Finding} -- This algorithm is the most in flux with no standards established by the community. However, the general approach first applies filters to the remaining image either a median or gaussian filter, followed by a simple or complex peak finding routine. After a Gaussian filter is applied, the beam is located through the maximum intensity pixel. More complex routines apply feature detection routines, novel masking approaches or even implement self-supervised vision models like DINOv3. Once the beam is located its intensity is integrated to determine the direct reflected beam intensity.
\item
  \textbf{Background Subtraction} -- A secondary background correction is then applied using a suitable region of the flattened detector far from the beam. A region the same size as the direct beam is integrated and its intensity is then subtracted. Usually, this final subtraction is nearly zero, but this process corrects for any remaining fluctuations or over corrections by the row-by-row correction.
\end{enumerate}

This routine produces a single intensity value for each incident angle, but the intensity is still only an absolute measurement that is only proportional to the number of electrons present within each physical pixel. To compute reflectance from this value, we must normalize by the measured incident intensity. This is captured by a set of frames at the start of each scan directly imaging the direct beam on the detector. To avoid damage, the higher-order suppressor is set to maximal attenuation and the exposure to the shutter limit, 0.001 s for the original shutter and 0.002 s for its replacement, with shot-to-shot variation quantified through repeated direct-beam frames (five for the ZnPc data of Chapter 4). The first reflectivity series is then captured under the same higher order suppressor and exposure time so that themeasured \(I_{0}\) frames can be used to directly normalize the series. This normalization is then propagated to each following series through a set of overlapping measurements at the end of the prior and start of the current series.

\begin{figure}[tbp]
  \centering
  \includegraphics[width=\linewidth,height=0.75\textheight,keepaspectratio]{fig2-09_reduction}
  \caption[Reduction of measured CCD frames into a stitched reflectivity curve, ZnPc at 285.0 eV in s polarization (Scan 85302).]{Reduction of measured CCD frames into a stitched reflectivity curve, ZnPc at 285.0 eV in s polarization (Scan 85302). a) A 2 × 2 binned frame at θ = 13.56° and 0.1 s on a logarithmic scale, with the specular region enlarged. Specular counts are summed in a 13 × 13 pixel box (solid), and the background is the mean of two equal boxes 28 pixels to either side (dashed). b) Background-subtracted count rate, counts divided by the exposure t plus the shutter offset t0 = 3.9 ms, for the four series and the direct beam I0 (star). Each series has its own exposure and higher-order suppressor position and repeats the last angles of the previous series. c) Stitched reflectivity with propagated uncertainties, mostly smaller than the markers, each series divided by I0 and by the running product of stitch factors Sn,n−1. Inset, R2/R1 at the overlap angles and S2,1 = 48.7 ± 2.0 with its uncertainty (red band).}
  \label{fig:3.5}
\end{figure}


Figure~\ref{fig:3.5} follows one measured scan through this reduction, a ZnPc film at 285.0 eV in s polarization collected on the CCD detector. On each frame the specular beam is summed in a fixed box and the background is estimated from two boxes of equal size on either side of it (Fig.~\ref{fig:3.5}a). Dividing the background-subtracted counts by the effective exposure \(t + t_{0}\) converts each frame to a count rate. The shutter offset \(t_{0}\) = 3.9 ms is nearly four times the 1 ms exposure of the first two series, the limit of the original shutter, so it cannot be neglected. The resulting count rates fall into four disjoint series (Fig.~\ref{fig:3.5}b). Within a series the rate follows the reflectivity, but at each series boundary the higher-order suppressor is withdrawn or the exposure lengthened, with the rate jumping by a factor that must be measured rather than taken from the nominal attenuation. The overlap angles provide that measurement. Each new series repeats the last few angles of the previous one, and the stitch factor \(S_{n,n - 1}\) is the slope of an orthogonal-distance regression of \(R_{n} = SR_{n - 1}\) through the origin on those angles (Fig.~\ref{fig:3.5}c, inset), as detailed below. Dividing the first series by the direct beam \(I_{0}\), and each later series by \(I_{0}\) and the running product of stitch factors, yields a single reflectivity curve spanning more than three decades in the first 0.10 Å\(^{- 1}\) (Fig.~\ref{fig:3.5}c). The stitch factors for this scan are \(S_{2,1} = 48.7 \pm 2.0\), \(S_{3,2} = 3.05 \pm 0.04\), and \(S_{4,3} = 1.13 \pm 0.03\). Their uncertainties enter every point of the later series as a common scale error, which a point-by-point uncertainty does not capture and which the error budget below identifies as a correlated term. The reduction shown here omits the dezinger and median filters and the beam-current normalization of the production code, pxr-reduce,\autocite{alsrsoxspxrreduce2026als} which implements the full routine described above.

\subsection{Stitching Details}\label{stitching-details}

Each series after the first opens by repeating the last angles of the series before it, and these overlap angles fix the relative scale of the two series (Fig.~\ref{fig:3.5}c, inset). Repeated frames at the first angle of a new series are excluded, and the remaining overlap angles \(k\) give paired reflectivities \(R_{n - 1,k}\) and \(R_{n,k}\), both normalized by the same \(I_{0}\) but recorded at different higher-order suppressor positions or exposure times. The stitch factor \(S_{n,n - 1}\) is the slope of the line \(R_{n,k} = S_{n,n - 1}R_{n - 1,k}\) through the origin, found by orthogonal distance regression\autocite{boggs1990orthogonal}, which minimizes the summed squared perpendicular distances of the overlap points from the line. Both series carry comparable counting noise at the overlap, and an ordinary least-squares fit, which assigns all of the scatter to the ordinate, would bias the slope low through the noise in the abscissa. The standard error of \(S_{n,n - 1}\) is taken from the covariance of the regression scaled by its residual variance, so it reflects the observed scatter of the overlap points rather than an assumed noise model. Series \(n\) is then placed on the scale of the first series by the running product of stitch factors, \(R_{n} = r_{n}/\left( I_{0}\prod_{m = 2}^{n}S_{m,m - 1} \right)\), where \(r_{n} = I_{n}/\left( t + t_{0} \right)\) is the background-subtracted count rate. The procedure follows pxr-reduce, with the reduction of Fig.~\ref{fig:3.5} performing the regression with the ODR module of SciPy.\autocite{virtanen2020scipy,alsrsoxspxrreduce2026als}

\subsection{Accounting for Measurement Uncertainty}\label{accounting-for-measurement-uncertainty}

Traditional modes of accounting for photon measurement errors originate with the standard counting (shot) noise that scales as the square root of the number of measured events. This assumes that the probability of an individual event does not depend on the number of events already measured. Under this approach the variance of a measured number of counts \emph{N} is \emph{N} itself.

This per-count shot noise propagates in quadrature through each step of the reduction routine above. The row-by-row and column-by-column background corrections each contribute an independent variance from the subtracted background estimate, and the subsequent background-region subtraction adds a further, independent term of the same form. Summed in quadrature, the total variance carried into each measured specular intensity is the sum of the counting variance, the row and column background variance, and the background-region variance, each following ordinary Poisson counting statistics before the detector\textquotesingle s own response is considered.

The uncertainty of each stitched point combines independent and shared terms. For \(R_{n}\) as above, the relative variance is \(\sigma_{R}^{2}/R^{2} = \sigma_{r}^{2}/r^{2} + \sigma_{I_{0}}^{2}/I_{0}^{2} + \sum_{m = 2}^{n}\sigma_{S,m}^{2}/S_{m,m - 1}^{2}\), where \(\sigma_{r}\) combines the counting variance of the specular box with the variances of the row, column, and region background estimates, and \(\sigma_{I_{0}}\) is the scatter of the repeated direct-beam frames. The first term changes from point to point. The last two are shared by every point of a series, so in a fit they act as a correlated scale error rather than as independent point uncertainties, although Fig.~\ref{fig:3.5} adds them in quadrature to each point as the reduction convention. For the scan of Fig.~\ref{fig:3.5} the stitch factors are \(S_{2,1} = 48.7 \pm 2.0\), \(S_{3,2} = 3.05 \pm 0.04\), and \(S_{4,3} = 1.13 \pm 0.03\), so the 4 percent uncertainty of the first stitch dominates the scale of every later series.

The detector records analog to digital units (ADU)\autocite{gann2012soft} a value that is proportional to the number of electrons generated within each pixel, with electron generation in turn proportional to the number of incident photons. The quantum efficiencies of the detector along with the gain set the photon-electron conversion rates, while the internal circuitry sets the conversion from electrons to recorded ADUs displayed on the detector. Notably the quantum efficiency and gain are independent of absolute intensity, yet the electron-ADU conversion depends on the number of electrons within each pixel well. The variance of the number of electrons generated per absorbed photon is also smaller than the Poisson value, by the Fano factor, which quantifies the sub-Poissonian suppression of variance in electron-hole pair generation. In practice this detector-level contribution is small next to the counting and background terms above and is neglected in the error budget below.

The dominant remaining source of uncertainty, however, is not detector noise but the ambiguity in separating the specular beam from off-specular diffuse scattering at high incident angle, where the two become difficult to resolve on a finite-sized detector. Chapter 4\textquotesingle s Methods addresses this directly with an empirical envelope term added in quadrature to the propagated counting uncertainty (Eq. 5), which dominates the reported uncertainty at high \emph{q}. The envelope is itself a model of the measurement error, and like the structural model it can be tested against the data and refined, for example by fitting its coefficients together with the structure.

\section{Analysis Methods}\label{analysis-methods}

Analysis of a reflectivity dataset couples a forward model, which computes the reflectivity of a parameterized film exactly, with an optimizer that adjusts the parameters to minimize the misfit to the data. A posterior estimate then converts the optimum into parameter uncertainties. Competing models are then compared by criteria that penalize added complexity. The sections below treat each stage in turn with the ZnPc film of Chapter 4 as the running example. Each stage is implemented in open-source packages, refloxide for the forward model and refnx for objectives, fitting, and sampling.\autocite{nelson2019textless,alsrsoxsrefloxide2026als}

\subsection{Simulating Reflectivity}\label{simulating-reflectivity}

All reflectivity in this dissertation is computed with refloxide, a 4 × 4 transfer matrix engine for polarized reflectivity from stratified media developed as part of this work.\autocite{alsrsoxsrefloxide2026als} For isotropic layers the reflectivity follows from the Parratt recursion at negligible cost, with well-tested implementations available.\autocite{parratt1954surface,nelson2019textless} For anisotropic layers each slab requires an eigendecomposition and a chain of 4 × 4 matrix products, and a resonant dataset repeats this for every energy, polarization, and angle measured, so the forward model dominates the cost of a fit. refloxide therefore places the calculation in a Rust core with thin Python bindings, offers an optional GPU path through wgpu, and plugs into refnx for modeling and fitting, so that a model spanning many energies remains practical inside an optimizer loop.

Near the critical angle the reflected amplitude is not small, and at every buried interface part of the field is reflected back into the film, reflected again at the interface above, and returned toward the detector with a phase set by the round trip (Fig.~\ref{fig:3.6}a). The transfer matrix method accounts for every such order exactly. Within each layer the field vector \(\Psi = \left( E_{x},H_{y},E_{y}, - H_{x} \right)^{\mathsf{T}}\) obeys \(\partial_{z}\Psi = i(\omega/c)\Delta\Psi\), where \(\Delta\) is the 4 × 4 Berreman matrix built from the local dielectric tensor\autocite{berreman1972optics,passler2017generalized,passler2019generalized}, whose eigenmodes define the interface matrices \(L_{i}\) and the propagation matrices \(P_{i}\). Their ordered product \(\Gamma_{N} = L_{1}P_{1}L_{2}P_{2}\cdots L_{N}P_{N}L_{N + 1}\) maps the field across the whole stack, and the reflectivities follow as \(R_{kl} = \left| r_{kl}\left( \Gamma_{N} \right) \right|^{2}\), valid because the incident medium is isotropic\autocite{yeh1980optics}. The isotropic limit reduces to the Parratt recursion,\autocite{parratt1954surface} and the anisotropic form is the one used for resonant reflectivity of oriented organic films\autocite{capelli2019quantitative,pasquali2014structural,pasquali2021analysis} and implemented in refloxide. Figure~\ref{fig:3.6}b to d compares this result with the kinematic approximation, which retains one reflection per interface, for three stacks at 285.1 eV. For the single ZnPc film of Chapter 4 (Fig.~\ref{fig:3.6}b) the kinematic curve exceeds unity below the critical angle and misplaces and deepens the first minima, although at higher \(q\) it reproduces the fringe spacing \(2\pi/t \approx 0.033\) Å\(^{- 1}\). For the edge-on over face-on bilayer (Fig.~\ref{fig:3.6}c) the agreement is closer, with a median deviation of 0.13 decades, but the kinematic model distorts the low-\(q\) fringes and drives the minimum near 0.18 Å\(^{- 1}\) below \(10^{- 8}\). In the six-period multilayer (Fig.~\ref{fig:3.6}d) both models place the Bragg peaks of the 80 Å period near 0.079 and 0.236 Å\(^{- 1}\), but below the first peak the kinematic model predicts sharp oscillations that multiple reflections damp and redistribute, with the median deviation reaching 0.31 decades. The kinematic errors are largest near the critical angle and at the deep minima, the regions that carry much of a fit\textquotesingle s sensitivity to density and interfacial structure (Fig.~\ref{fig:3.8}),so every simulation in this dissertation uses the transfer matrix.

\begin{figure}[tbp]
  \centering
  \includegraphics[width=\linewidth,height=0.75\textheight,keepaspectratio]{fig2-11_transfer-matrix}
  \caption[Multiple reflections make the exact solution of the layer stack necessary.]{Multiple reflections make the exact solution of the layer stack necessary. a) Every interface of a stack with edge-on (layer 1) and face-on (layer 2) molecules reflects and transmits, and all reflected orders interfere at the detector. The field vector Ψ obeys ∂zΨ = i(ω/c)ΔΨ in each uniaxial layer, and interface matrices Li and propagation matrices Pi compose to the stack matrix ΓN that gives the reflection coefficients. b to d) s-polarized reflectivity at 285.1 eV from the transfer matrix (red) and the kinematic approximation (black, dashed) for the Chapter 4 single film (b), a 100 Å edge-on layer over a 90 Å face-on layer (c), and six periods of 40 Å edge-on and 40 Å face-on layers (d). The median deviation between the two is 0.23, 0.13, and 0.31 decades for b, c, and d.}
  \label{fig:3.6}
\end{figure}


\begin{figure}[tbp]
  \centering
  \includegraphics[width=\linewidth,height=0.75\textheight,keepaspectratio]{fig2-12_interfaces}
  \caption[Three interface models at the same roughness σ = 20 Å for a 192.7 Å ZnPc film on Si at 283.7 eV.]{Three interface models at the same roughness σ = 20 Å for a 192.7 Å ZnPc film on Si at 283.7 eV. a) Real-space interfaces drawn at equal lateral and vertical scale, with the wavelength λ = 43.7 Å and 2σ (red) marked. Névot-Croce describes fine roughness with a short in-plane correlation length ξ, Debye-Waller describes broad facets with a long ξ, and the graded interface describes interdiffusion without in-plane structure. b) Depth profiles of δxx as each model computes them. Névot-Croce and Debye-Waller compute the sharp profile (red), shown with the error-function profile usually drawn for them (black, dashed), while the graded interface slices the error-function profile into 1 Å layers and computes it directly (red), with the sharp profile for reference (gray, dashed).}
  \label{fig:3.7}
\end{figure}


Real interfaces are neither flat nor sharp. The three models in common use differ in what they assume about an interface and in where that assumption enters the calculation (Fig.~\ref{fig:3.7}). Névot-Croce and Debye-Waller roughness both keep the computed depth profile sharp and multiply each interface Fresnel coefficient by a damping factor, \(\exp\left( - 2k_{z,i}k_{z,i + 1}\sigma^{2} \right)\) for Névot-Croce\autocite{nevot1980caracterisation} and \(\exp\left( - 2k_{z,i}^{2}\sigma^{2} \right)\) for Debye-Waller\autocite{deboer1996x,holy1993x}. The error-function profile usually drawn for these models (Fig.~\ref{fig:3.7}b, dashed) is a display convention that the calculation never uses, so the two models share one depth profile and differ only in the reflectivity they predict. The damping factors encode different interfaces. Névot-Croce assumes fine roughness with an in-plane correlation length \(\xi\) short compared with \(\xi_{diff} = \lambda/(1 - cos\theta)\), whereas Debye-Waller assumes broad facets with \(\xi\) long compared with \(\xi_{diff}\) (Fig.~\ref{fig:3.7}a). At 283.7 eV, \(\xi_{diff}\) falls from about \(10^{4}\) Å at \(\theta\) = 5° to about 300 Å at \(\theta\) = 30°, so the appropriate factor can change within a single scan and depends on lateral structure that specular reflectivity does not measure. The graded interface makes no claim about lateral structure. It slices the error-function profile into 1 Å layers and propagates the field through each one, so the depth profile is the computed one (Fig.~\ref{fig:3.7}b, right). This is the appropriate description of interdiffusion and of real composition or orientation gradients, serving as the reference for the multiplicative models. At \(\sigma\) = 20 Å the s-polarized reflectivity from the Névot-Croce and Debye-Waller factors departs from the graded result by up to 28 and 51 percent below \(q\) = 0.15 Å\(^{- 1}\), so at large roughness the factors are approximations to be tested against the graded model, consistent with numerical comparisons of these roughness treatments for X-ray and extreme-ultraviolet reflectance.\autocite{esashi2021influence} The refloxide selection guide states these regimes in terms of \(\sigma k_{z}\) and \(\xi/\xi_{diff}\).\autocite{alsrsoxsrefloxide2026als} In Chapter 4 the slab model applies Névot-Croce roughness at each interface, while the depth gradient is a physically motivated multi-slice profile of molecular tilt and density, computed slice by slice in the same way as the graded interface.

The graded interface is one case of a general multi-slice description, in which any depth profile of the optical tensor is discretized into slabs thin compared with \(\pi/q_{z,max}\) and each slab is assigned its own tensor. A composition gradient varies the density of each component across the slices, and an orientation gradient varies the tilt \(\gamma\) entering the film tensor of Chapter 2, so that \(n_{xx}^{f}(z)\) and \(n_{zz}^{f}(z)\) change independently with depth. The exponential profiles of the orientation gradient section enter the Chapter 4 model in this way. The slice thickness should be chosen so that halving it changes the computed reflectivity by less than the measurement uncertainty.

\subsection{Fitting Polarized Reflectivity}\label{fitting-polarized-reflectivity}

The goal of the resonant scattering experiment is to uncover physical structures and size scales along with the molecular states within these structures, but there are no general direct inversion algorithms to uncover both properties. Because scattering intensity is measured in the far field regime, it arises from the Fourier transform of the near field projection, with much of the near-field information lost because only intensity is measured. As a direct consequence, the only way to quantitatively uncover both properties follows an inverse problem. A model and parameter set are chosen and used to simulate the experimental scattering, after which the parameters are refined to fit the data. But this requires the user to choose a good model that can describe the underlying structures and molecular states along with a good choice in parameterization that sufficiently removes covariance and degeneracies. Yet even for good choices, there is no guarantee of solution uniqueness, since multiple models and parameterization choices can result in statistically identical simulated scattering intensities.

Because of this, it is the job of the researcher to identify the most likely model based on physical and chemical properties and other experimental results. Prior chemical knowledge, such as the expected film stoichiometry or known crystallographic phases, together with auxiliary measurements like GIXRD or ellipsometry, constrains this choice before any fitting begins. These constraints enter a fit as fixed parameters, such as the optical constants of the substrate, as bounds that exclude unphysical densities or thicknesses, and as joint fits in which one structure must describe several energies and polarizations at once. The joint fit is the strongest of these, and it is the constraint that the physically constructed optical tensors of Chapters 4 and 5 make available.

\begin{figure}[tbp]
  \centering
  \includegraphics[width=\linewidth,height=0.75\textheight,keepaspectratio]{fig2-13_objective-transforms}
  % TODO caption contains equations that were flattened to plain text in the Word export; retypeset them.
  \caption[Objective transforms change which points control a reflectivity fit.]{Objective transforms change which points control a reflectivity fit. A simulated s-polarized dataset for the Chapter 4 ZnPc film at 285.1 eV, with noise σ=aR+bq (points) and the model (red), shown as a) R, b) log10R, and c) Rq4. Below each panel is the statistical leverage of every point in an unweighted least-squares fit of film thickness, roughness, tilt, and density, relative to a uniform share (dashed line). Neff/N is the participation ratio, the fraction of points that effectively carry the fit.}
  \label{fig:3.8}
\end{figure}


\subsection{Objective Transforms - linearization}\label{objective-transforms---linearization}

An unweighted least-squares objective counts every point equally in whatever representation the data are fit, so the representation decides which points control the fit. Figure~\ref{fig:3.8} shows one simulated dataset, the Chapter 4 film at 285.1 eV with noise \(\sigma = aR + bq\), as \(R\), \(\log_{10}R\), and \(Rq^{4}\) (Fig.~\ref{fig:3.8}a to c). Below each panel is the statistical leverage of every point, the diagonal of the hat matrix \(\mathbf{J}\left( \mathbf{J}^{\mathsf{T}}\mathbf{J} \right)^{- 1}\mathbf{J}^{\mathsf{T}}\) for film thickness, roughness, tilt, and density, relative to a uniform share (dashed line). Leverage measures how strongly the fitted model is pulled toward each point, with its participation ratio \(N_{eff}/N\) giving the fraction of the dataset that effectively carries the fit. On a linear scale the critical-angle region dominates, with the central half of the leverage between \(q\) = 0.014 and 0.046 Å\(^{- 1}\) (\(N_{eff}/N\) = 0.20). The logarithm spreads weight across the decades but concentrates it at the fringe minima, where a small shift in fringe position changes \(\log R\) by orders of magnitude. The deep minimum near 0.135 Å\(^{- 1}\) alone carries 93 times a uniform share (\(N_{eff}/N\) = 0.04). Multiplying by \(q^{4}\) removes the Fresnel decay, so the leverage stays close to the uniform share across most of the range, never exceeding 2.2 times it, with three quarters of the points contributing effectively (\(N_{eff}/N\) = 0.75). With correctly propagated uncertainties the choice matters little, since \(\sigma_{\log R} = \sigma_{R}/\left( R\ln 10 \right)\) and to first order every representation yields the same weighted \(\chi^{2}\). It matters when uncertainties are absent, underestimated, or replaced by uniform weights. Because each transform also changes the implied distribution of the noise, the representation is part of the statistical model and is reported with it. Both transforms are available as objective transforms in refnx\autocite{nelson2019textless} and are used in this dissertation. The fits of Chapter 4 and the ZnPc comparison of Chapter 5 use \(\log R\), while the Alq\(_{3}\) fits of Chapter 5 use \(Rq^{4}\).

\subsection{Optimizers}\label{optimizers}

Reflectivity objectives combine three difficulties, isolate in Fig.~\ref{fig:3.9} on a two-parameter problem. The Kiessig fringes make the objective sharp and periodic, since a few ångströms of thickness move every fringe. The total optical path of the film is well determined while its division between layers is not, which correlates the parameters into long, narrow, tilted valleys. Fringe aliasing repeats those valleys wherever the model fringes fall on the data fringes shifted by whole periods, each valley a local minimum. The surface in Fig.~\ref{fig:3.9} is built from simulated s- and p-polarized data for the Chapter 4 film with its buriedinterfacial layer, with every optimizer starts from the same point, 17 Å and 15 Å from the true thicknesses. Levenberg-Marquardt, which interpolates between Gauss-Newton and gradient descent,\autocite{marquardt1963algorithm,levenberg1944method} and L-BFGS-B, which builds a limited-memory quasi-Newton Hessian within box bounds,\autocite{byrd1995limited} both descend into the nearest valley and stop at (210.4 Å, 14.1 Å) with \(\log_{10}\chi_{\nu}^{2} = 1.52\), after 69 and 48 evaluations (Fig.~\ref{fig:3.9}a,b). Differential evolution needs no starting point.\autocite{storn1997differential} It evolves a population across the bounded space by mutation and crossover, and by generation 39 the population has collapsed onto the global minimum at (178.0 Å, 15.0 Å), \(\log_{10}\chi_{\nu}^{2} = - 0.58\), at a cost of 984 evaluations (Fig.~\ref{fig:3.9}c). A random-walk Metropolis chain\autocite{metropolis1953equation} also reaches the minimum from this start, but in the sharp valley it accepts only 1 percent of its proposals, so 4000 steps yield few independent samples (Fig.~\ref{fig:3.9}d). We expect a chain started in another valley to cross into this one rarely, which is the mixing problem of naive MCMC on reflectivity objectives.\autocite{machida2024development} Local methods are therefore fast and precise only inside the correct basin, global methods find that basin at 14 to 20 times the cost. Sampling methods are best used to characterize a basin once it has been found. The Chapter 4 workflow follows this order, with differential evolution as the global stage and L-BFGS-B as the local polish.

\begin{figure}[tbp]
  \centering
  \includegraphics[width=\linewidth,height=0.75\textheight,keepaspectratio]{fig2-15_uncertainty}
  \caption[From the global optimum to parameter uncertainty on the Fig.]{From the global optimum to parameter uncertainty on the Fig.~\ref{fig:3.9} objective. a) Differential evolution with the population at generations 0 (gray), 4 (light red), and 8 (red), reaching (178.2 Å, 14.9 Å) after 270 evaluations. The red box marks the basin enlarged in b and c. b) The log10χ2ν surface over that basin (0.83 Å by 0.46 Å) with the differential-evolution best point (star) and the local polish by L-BFGS-B (red, 27 evaluations) and Nelder-Mead (black, 58 evaluations), both converging to (178.00 Å, 14.99 Å). c) Affine-invariant ensemble MCMC in the same window, 32 walkers for 3000 steps after 1000 steps of burn-in, with an acceptance fraction of 0.72. d) Posterior samples with marginal distributions and 1σ and 2σ contours from the MCMC covariance (black, dashed), the finite-difference Hessian at the minimum (red), and the L-BFGS-B inverse Hessian (red, dotted).}
  \label{fig:3.10}
\end{figure}


\begin{figure}[tbp]
  \centering
  \includegraphics[width=\linewidth,height=0.75\textheight,keepaspectratio]{fig2-14_optimizers}
  \caption[Four optimizers on the same sharp, correlated, and nonlinear reflectivity objective.]{Four optimizers on the same sharp, correlated, and nonlinear reflectivity objective. The surface is log10χ2ν of simulated s- and p-polarized reflectivity at 283.7 eV for the Chapter 4 film divided into a bulk layer of thickness tbulk and a buried interfacial layer of thickness tint, fit on log R, with a single deep minimum (star) at tbulk = 178 Å and tint = 15 Å. Every run starts from the same point (black circle, 195 Å, 30 Å). a, b) Levenberg-Marquardt (69 evaluations) and L-BFGS-B (48 evaluations) descend into the nearest valley and stop at the same local minimum, (210.4 Å, 14.1 Å) with log10χ2ν = 1.52. c) Differential evolution population at generations 0 (gray), 4 (light red), and 39 (red), which converges to the global minimum (178.0 Å, 15.0 Å) with log10χ2ν = −0.58 after 984 evaluations. d) Random-walk Metropolis chain of 4000 steps at the true likelihood, which reaches the minimum from this start but accepts 1 percent of its proposals in the sharp valley.}
  \label{fig:3.9}
\end{figure}


\subsection{Posterior Estimation}\label{posterior-estimation}

Figure~\ref{fig:3.10} completes the workflow on the same objective by converting the global optimum into a parameter uncertainty. Differential evolution locates the basin (Fig.~\ref{fig:3.10}a), but its population carries no curvature information, so it returns a best point without an error estimate. A local optimizer started from that point refines it, and L-BFGS-B and the derivative-free Nelder-Mead method converge to the same minimum at (178.00 Å, 14.99 Å) in 27 and 58 evaluations (Fig.~\ref{fig:3.10}b). Within this basin the objective is close to quadratic, the regime in which affine-invariant ensemble MCMC,\autocite{goodman2010ensemble} as implemented in emcee and refnx\autocite{nelson2019textless}, samples efficiently. Thirty-two walkers accept 72 percent of their proposals (Fig.~\ref{fig:3.10}c), against 1 percent for the random-walk chain on the full surface of Fig.~\ref{fig:3.9}d. The samples give \(\sigma\left( t_{bulk} \right)\) = 0.099 Å and \(\sigma\left( t_{int} \right)\) = 0.050 Å with a correlation of \(- 0.63\) (Fig.~\ref{fig:3.10}d), and the inverse of the finite-difference Hessian at the minimum reproduces these values to within 0.001 Å, as expected for a quadratic basin. The inverse Hessian reported by L-BFGS-B does not. It overstates \(\sigma\left( t_{bulk} \right)\) at 0.150 Å and understates the correlation at \(- 0.40\), because L-BFGS-B builds its Hessian estimate from the few steps it takes, which are few for a run started near the minimum. The internal estimate of a quasi-Newton optimizer is therefore a starting point to be checked against a sampled posterior or a Hessian computed at the minimum. Where the posterior has separated modes, as on the full surface of Fig.~\ref{fig:3.9}, nested sampling as implemented in dynesty\autocite{speagle2020dynesty} is the appropriate extension, since it explores separated modes and returns the evidence used for model comparison in the section that follows. Other approaches such as neural posterior estimation and Bayesian optimization are other approaches to resolve this.\autocite{starostin2025fast,machida2024development}

\subsection{Model Selection}\label{model-selection}

The default statistical test for evaluating model performance considers only the closeness between a model and the data it aims to fit. This is best characterized by the reduced \(\chi^{2}\) statistic that measures this directly. However, this assumes Gaussian, locally linear errors and penalizes nothing for added model complexity, so it cannot by itself discriminate between parameterizations of different dimensionality when the noise is closer to Poissonian and the forward model is nonlinear, as is generically the case for reflectivity.

Other evaluation criteria instead compare models directly against each other. The Bayesian and Akaike information criteria (BIC, AIC) approximate, respectively, the weight of evidence for a model and its expected predictive performance on new data, penalizing additional free parameters so a better fit is only rewarded if it earns its added complexity. Both remain asymptotic approximations, however, and the most direct test of generalization is a validation hold-out, fitting to a subset of the measured energies or angles and scoring the fit against the remainder, which catches overfitting that an information criterion can only estimate.

\subsection{Bayesian Optimization and Amortized Inference}\label{bayesian-optimization-and-amortized-inference}

Bayesian optimization and amortized inference reduce the number of forward evaluations a fit requires rather than improving the search itself. Bayesian optimization fits a Gaussian-process surrogate to the objective and chooses each next evaluation by an acquisition function that balances the predicted objective against its uncertainty.\autocite{noack2021gaussian} It pays off when one forward evaluation is expensive, as for a finite-element solve or a model spanning many energies, and when the number of parameters is small. It is the same machinery proposed for measurement selection in the future-work section. Prior-amortized neural posterior estimation instead trains a normalizing-flow posterior once on simulated reflectivity over a family of priors.\autocite{starostin2025fast} Inference on a new dataset then takes seconds, recovers modes that maximum-likelihood fitting and local MCMC miss, and can be refined by importance sampling or MCMC. Its published demonstrations use isotropic layer models at a single energy. Applying it to the uniaxial models of this work would require training sets simulated with the 4 × 4 transfer matrix at every measured energy and both polarizations, a network conditioned on the energy-dependent tensor of each component or on the parameters that generate it, and a forward model fast enough to make those training sets affordable, which is one motivation for the performance of refloxide.

\section{\texorpdfstring{Future Work }{Future Work }}\label{figmarker19x-figmarker20x-future-work}

\begin{figure}[tbp]
  \centering
  \includegraphics[width=\linewidth,height=0.75\textheight,keepaspectratio]{fig3-11_beamline-ui}
  \caption[Future Beamline 11.0.1.2 beamline user interface (left, light mode, right dark mode).]{Future Beamline 11.0.1.2 beamline user interface (left, light mode, right dark mode). Integrates pan automation, data access and analysis, and a standard control integration to directly drive the beamline.}
  \label{fig:3.11}
\end{figure}


The analysis outlined here is implemented in a Python package co-developed with the beamline staff at the Advanced Light Source.\autocite{alsrsoxspxrreduce2026als} The next step builds the entire measurement process directly into the control systems and drives for the instrument to facilitate the real time reduction and analysis. Importantly, the facility and instruments are in the process of upgrading from a LabVIEW-based control system to a newer expressive system called EPICS. These new EPICS drivers are then integrated directly with Bluesky a common control language used by most other synchrotron facilities and CERN. Bluesky makes it possible to safely write python scripts to control the beamline optics and motors, unlocking the potential for feedback loops and programmatic interaction. Bluesky also integrates with more modern features such as web applications allowing for more standardized UI experience and centralized data storage to unify the experience of collecting and working with data across user facilities. Beamline 11.0.1.2 is currently in the process of upgrading both to use EPICS and Bluesky and a UI/UX refresh that will utilize modern web features such as common authentication and finer-grained, per-user access control (Fig.~\ref{fig:3.11}).

This unlocks direct integration into decision making algorithms such as gaussian process (GP) routines so that new configurations can be algorithmically selected based on the uncertainties in the measured signal, completely removing the need for pre-defined human input such as defining the energy measurement grid. Using a gaussian process has been highly effective in reducing the number of measurements needed to fully capture experimental data, outcompeting uniform grid searches, human decision making and random grids, while also allowing for domain specific and model specific information injection.\autocite{du2025demonstration,alghalayini2025machine} The general idea is to use direct feedback algorithms to automate both the alignment routines and the single condition measurements, while the GP algorithm is used to set a stopping condition on the measurement grid, asserting that each measurement condition (polarization and energy) is fully specified. The next approach then directly integrates the materials film structure into the decision-making process, from the fit performance, using Bayesian optimization (BO) the stopping condition is no longer when the reflectivity is fully specified but when the film's depth profile is fully determined. In this final approach the user of the facility no longer needs to perform the primary reflectivity analysis on their films after their experiment is done, and instead they return home with the depth profile itself allowing their work to focus on uncovering material properties, structure property relationships in functionalized applications, and revealing new phenomena within novel materials.

Figure~\ref{fig:3.12} assembles the procedures of this chapter into the decision loop planned for beamline 11.0.1.2 once its EPICS and Bluesky controls are in place,\autocite{allan2019blueskytextquotesingle} following the cross-facility approach to autonomous scattering.\autocite{koepp2026toward} A Gaussian process holds the current state of knowledge\autocite{noack2021gaussian} and chooses eitherthe next measurement condition, a combination of polarization, photon energy, and angular range, or the next sample. Both branches pass through an alignment checker that encodes the rules of the alignment section. A change of polarization changes the heat load on the optics, so the loop waits 15 min before checking, with the sample realigned with the scans of Fig.~\ref{fig:3.4} only when it fails a check of the reflected beam at 1° and 2°. A condition is then measured one series at a time. Before each series the incident intensity is raised through the higher-order suppressor and the exposure time, and the series runs until the reflected beam falls below the noise floor, which reproduces the series structure of Fig.~\ref{fig:3.5} without a predefined list of angles. When a condition is complete, control returns to the Gaussian process, whose stopping rule asserts that every polarization and energy has been specified. Replacing that rule with Bayesian optimization on the fit itself would end the measurement when the depth profile, rather than the reflectivity, is determined. The loop remains under development and is outlined here to set the direction rather than to report a result.

\begin{figure}[tbp]
  \centering
  \includegraphics[width=\linewidth,height=0.75\textheight,keepaspectratio]{fig2-16_automation}
  \caption[Decision loop for autonomous reflectivity measurement.]{Decision loop for autonomous reflectivity measurement. A Gaussian process chooses the next condition (polarization, photon energy, and angle range) or the next sample, and a new sample is located before both branches pass to the alignment checker. After a polarization change the loop waits 15 min, and the sample is realigned whenever it fails the check at 1° and 2°. Each condition is measured one series at a time, with the incident intensity optimized for each series and data collected until the beam is lost, and a completed condition (red) returns control to the Gaussian process.}
  \label{fig:3.12}
\end{figure}


\section{Summary}\label{summary}

This chapter has assembled the physical and practical basis on which every later result rests. The detector records only far-field intensity, so the phase and near-field structure of the scattered wave are lost, and every result in this dissertation is therefore obtained by forward modeling the measurement and fitting it. The optical response entering that model is the tensor index \(n = 1 - \delta + i\beta\), with \(\beta\) taken from the mass absorption and \(\delta\) from its Kramers-Kronig transform, so the bare-atom extrapolation of \(\beta\) outside the measured window enters \(\delta\) at every energy. Each resonance contributes to this tensor through its transition dipole, and averaging the uniaxial molecular tensor over the azimuth gives the film tensor as a function of the tilt \(\gamma\) and the density \(\rho\), the two quantities refined in Chapters 4 and 5. The Ginzburg-Landau treatment adds a physical form for how these quantities relax toward the bulk, the exponential depth profiles used for the interfacial gradients of Chapter 4.

The experimental and analysis sections carry this framework into practice. In the measurement, the angular zero is set by a forward-model fit of the rocking scan, because a finite beam and photodiode bias the conventional half-maximum alignment, and the stitch factors that join attenuated series carry a correlated scale uncertainty that a point-by-point error bar does not capture. In the analysis, reflectivity is computed exactly by the transfer matrix, roughness is treated with Névot-Croce factors checked against explicitly graded interfaces, and fits are made on \(\log R\) or \(Rq^{4}\) with the transform reported as part of the noise model, since the transform decides which points control the fit. The global optimum is located by differential evolution, refined with L-BFGS-B, and assigned uncertainties by ensemble MCMC, a sequence suited to the sharp, correlated, and multimodal objectives typical of reflectivity. The chapter closes with the autonomous measurement loop toward which beamline 11.0.1.2 is moving, in which the same forward model decides what to measure next.
